% ELT-1031 Lab NN -- <lab title>
% Build: latexmk -lualatex phelan-labNN.tex
\documentclass{phelanlab}

\graphicspath{{figures/}}

% ---------- Shorthands ----------
% Circuit quantities: upright with a real subscript in text, italic in maths.
\newquantity{\Vs}{V}{S}
\newquantity{\Vth}{V}{TH}
\newquantity{\Rth}{R}{TH}
\newquantity{\Req}{R}{EQ}

\title{Lab Title from the Handout}
\author{Leonard H Phelan IV}
\date{2026-10-12}            % lab date, YYYY-MM-DD
\labnumber{NN}
% \course{ELT-1031}            % default
% \instructor{prof. Ashworth}  % default
\footertext{Phelan\enspace·\enspace ELT-1031\enspace·\enspace Lab NN}

\begin{document}
\maketitle

% =====================================================================
\section*{Objective}

One paragraph on what the lab sets out to show. Numbers and units go through
\verb|\qty|, which keeps them together on one line: \qty{10.01}{\volt},
\qty{2.3}{\percent}, \qtyrange{0.7}{2.7}{\percent}. Labels and their numbers
take a tie: Part~I, Table~\ref{tab:example}.

\section*{Equipment}

Siglent SPD3303X-E DC power supply; Siglent SDM3055A bench DMM; NI Multisim.

% =====================================================================
\section{Part I: First Part Title}

In text, \Vth{} and \Rth{} are upright with real subscripts. In maths the
letter is italic:
\begin{equation*}
  \Rth = \qty{100}{\ohm} + (\qty{300}{\ohm} \parallel \qty{200}{\ohm})
       = \qty{220}{\ohm}
\end{equation*}

\subsection{Measurements}

% Column widths are weights: X[2] is twice as wide as X[1].
\begin{table}
  \begin{labtable}{X[1.4]X[1]X[1]X[1]}
    Quantity & Calculated & Measured & err \\
    \Vs      & \qty{10.00}{\volt} & \qty{10.01}{\volt} & \qty{+0.10}{\percent} \\
    \Vth     & \qty{6.000}{\volt} & \qty{6.02}{\volt}  & \qty{+0.33}{\percent} \\
    \Rth     & \qty{220.0}{\ohm}  & \pending           & \pending \\
  \end{labtable}
  \caption{The caption goes in the margin, level with the bottom of the table.}
  \label{tab:example}
\end{table}

% \fig{width as a fraction of the text width}{file}{caption}{label}
% \fig{0.6}{p1-ms-original}{Original circuit in Multisim.}{fig:p1-ms-original}

\subsection{Results}

What the numbers show.

% =====================================================================
\section*{Questions}

\begin{question}
  Restate the handout question here.
\end{question}
\begin{answer}
  The answer, indented under the question and kept on one page.
\end{answer}

\begin{question}[Q2b]
  A question with its own label, for when the handout numbering isn't 1, 2, 3.
\end{question}
\begin{answer}
  Another answer.
\end{answer}

% =====================================================================
\section*{Errors and Difficulties}

What went wrong and what was done about it.

\end{document}
